%%%%%%%%%%%%%%%%%%%%%%%%%%%%%%%%%%%%%%%%%%%%%%%%%%%%%%%%%%%%%%%%%%%%%%
% Câu chưa có mã dạng — gắn %[..] rồi chuyển vào đúng file dạng
%%%%%%%%%%%%%%%%%%%%%%%%%%%%%%%%%%%%%%%%%%%%%%%%%%%%%%%%%%%%%%%%%%%%%%


% nguồn: De_AI_Toan9_cuoi_hoc_ki_i_1 (nhập 18/11/2026)
\begin{ex}%[ma:03f6ede0]
	Phương trình nào sau đây là phương trình bậc nhất hai ẩn?
	\choice
	{\True $2x + y = 3$}
	{$0x + 0y = 1$}
	{$x^2 + y = 0$}
	{$x + y + z = 1$}
	\loigiai{
		Phương trình bậc nhất hai ẩn có dạng $ax + by = c$, trong đó $a, b$ không đồng thời bằng $0$. Do đó $2x + y = 3$ là phương trình bậc nhất hai ẩn.
	}
\end{ex}

% nguồn: De_AI_Toan9_cuoi_hoc_ki_i_1 (nhập 18/11/2026)
\begin{ex}%[ma:1ee4d3d9]
	Hệ phương trình nào sau đây nhận cặp số $(1; -1)$ làm nghiệm?
	\choice
	{$\begin{cases} x - y = 0 \\ x + y = 2 \end{cases}$}
	{\True $\begin{cases} x + y = 0 \\ 2x - y = 3 \end{cases}$}
	{$\begin{cases} 2x + y = 1 \\ x - y = 1 \end{cases}$}
	{$\begin{cases} x + y = 1 \\ 2x + y = 1 \end{cases}$}
	\loigiai{
		Thay $x = 1, y = -1$ vào hệ phương án B ta được: $1 + (-1) = 0$ (đúng) và $2(1) - (-1) = 3$ (đúng).
	}
\end{ex}

% nguồn: De_AI_Toan9_cuoi_hoc_ki_i_1 (nhập 18/11/2026)
\begin{ex}%[ma:7e955b95]
	Tập nghiệm của phương trình $(x - 2)(2x + 1) = 0$ là:
	\choice
	{$\{2\}$}
	{$\{-0{,}5\}$}
	{\True $\{2; -0{,}5\}$}
	{$\{-2; 0{,}5\}$}
	\loigiai{
		Ta có $(x - 2)(2x + 1) = 0 \Leftrightarrow x - 2 = 0$ hoặc $2x + 1 = 0 \Leftrightarrow x = 2$ hoặc $x = -0{,}5$.
	}
\end{ex}

% nguồn: De_AI_Toan9_cuoi_hoc_ki_i_1 (nhập 18/11/2026)
\begin{ex}%[ma:40f8e356]
	Cho $a > b$. Bất đẳng thức nào sau đây đúng?
	\choice
	{$a - 2 < b - 2$}
	{$-2a > -2b$}
	{$2a + 1 < 2b + 1$}
	{\True $a + 3 > b + 3$}
	\loigiai{
		Cộng hai vế của bất đẳng thức $a > b$ với $3$, ta được $a + 3 > b + 3$.
	}
\end{ex}

% nguồn: De_AI_Toan9_cuoi_hoc_ki_i_1 (nhập 18/11/2026)
\begin{ex}%[ma:ff409f0d]
	Nghiệm của bất phương trình $-2x + 4 > 0$ là:
	\choice
	{\True $x < 2$}
	{$x > 2$}
	{$x < -2$}
	{$x > -2$}
	\loigiai{
		Ta có $-2x + 4 > 0 \Leftrightarrow -2x > -4 \Leftrightarrow x < 2$ (chia hai vế cho $-2$ và đổi chiều).
	}
\end{ex}

% nguồn: De_AI_Toan9_cuoi_hoc_ki_i_1 (nhập 18/11/2026)
\begin{ex}%[ma:27c5f0c0]
	Biểu thức $\sqrt{x - 3}$ xác định khi và chỉ khi:
	\choice
	{$x \le 3$}
	{\True $x \ge 3$}
	{$x > 3$}
	{$x \ne 3$}
	\loigiai{
		Biểu thức căn thức bậc hai xác định khi biểu thức dưới dấu căn không âm, tức là $x - 3 \ge 0 \Leftrightarrow x \ge 3$.
	}
\end{ex}

% nguồn: De_AI_Toan9_cuoi_hoc_ki_i_1 (nhập 18/11/2026)
\begin{ex}%[ma:ecbb9eb0]
	Giá trị của biểu thức $\sqrt{12} - 2\sqrt{3}$ bằng:
	\choice
	{$2\sqrt{3}$}
	{$-2\sqrt{3}$}
	{\True $0$}
	{$4\sqrt{3}$}
	\loigiai{
		Ta có $\sqrt{12} - 2\sqrt{3} = \sqrt{4 \cdot 3} - 2\sqrt{3} = 2\sqrt{3} - 2\sqrt{3} = 0$.
	}
\end{ex}

% nguồn: De_AI_Toan9_cuoi_hoc_ki_i_1 (nhập 18/11/2026)
\begin{ex}%[ma:4f2cfafb]
	Căn bậc ba của $-8$ là:
	\choice
	{$2$}
	{$\pm 2$}
	{$-8$}
	{\True $-2$}
	\loigiai{
		Vì $(-2)^3 = -8$ nên $\sqrt[3]{-8} = -2$. Căn bậc ba của một số âm là một số âm.
	}
\end{ex}

% nguồn: De_AI_Toan9_cuoi_hoc_ki_i_1 (nhập 18/11/2026)
\begin{ex}%[ma:d1465084]
	Cho tam giác $ABC$ vuông tại $A$. Tỉ số lượng giác $\sin B$ bằng:
	\choice
	{\True $\frac{AC}{BC}$}
	{$\frac{AB}{BC}$}
	{$\frac{AC}{AB}$}
	{$\frac{AB}{AC}$}
	\loigiai{
		Trong tam giác vuông, $\sin = \frac{\text{cạnh đối}}{\text{cạnh huyền}}$. Với góc $B$, cạnh đối là $AC$, cạnh huyền là $BC$ nên $\sin B = \frac{AC}{BC}$.
	}
\end{ex}

% nguồn: De_AI_Toan9_cuoi_hoc_ki_i_1 (nhập 18/11/2026)
\begin{ex}%[ma:2e445d7f]
	Cho tam giác $ABC$ vuông tại $A$, đường cao $AH$. Biết $BH = 2\text{ cm}$ và $CH = 8\text{ cm}$. Độ dài $AH$ bằng:
	\choice
	{$16\text{ cm}$}
	{\True $4\text{ cm}$}
	{$8\text{ cm}$}
	{$10\text{ cm}$}
	\loigiai{
		Áp dụng hệ thức lượng trong tam giác vuông: $AH^2 = BH \cdot CH = 2 \cdot 8 = 16$. Suy ra $AH = \sqrt{16} = 4\text{ (cm)}$.
	}
\end{ex}

% nguồn: De_AI_Toan9_cuoi_hoc_ki_i_1 (nhập 18/11/2026)
\begin{ex}%[ma:6ce2ea34]
	Một đường tròn có bao nhiêu trục đối xứng?
	\choice
	{$1$}
	{$2$}
	{\True Vô số}
	{Không có}
	\loigiai{
		Mỗi đường thẳng đi qua tâm đều là một trục đối xứng của đường tròn. Vì có vô số đường thẳng đi qua tâm nên đường tròn có vô số trục đối xứng.
	}
\end{ex}

% nguồn: De_AI_Toan9_cuoi_hoc_ki_i_1 (nhập 18/11/2026)
\begin{ex}%[ma:4722962e]
	Cho đường tròn $(O; 5\text{ cm})$ và đường thẳng $a$. Khoảng cách từ tâm $O$ đến đường thẳng $a$ là $3\text{ cm}$. Số điểm chung của đường thẳng $a$ và đường tròn $(O)$ là:
	\choice
	{$0$}
	{$1$}
	{$3$}
	{\True $2$}
	\loigiai{
		Gọi $d$ là khoảng cách từ $O$ đến $a$, $R$ là bán kính. Vì $d = 3 < 5 = R$ nên đường thẳng cắt đường tròn tại $2$ điểm phân biệt.
	}
\end{ex}

% nguồn: De_AI_Toan9_cuoi_hoc_ki_i_1 (nhập 18/11/2026)
\begin{ex}%[ma:494ca6ad]
	Một bạn học sinh mua $3$ quyển vở và $2$ chiếc bút hết $36\,000$ đồng. Nếu mua $2$ quyển vở và $4$ chiếc bút loại đó thì hết $40\,000$ đồng. Gọi $x, y$ (nghìn đồng) lần lượt là giá tiền của một quyển vở và một chiếc bút. Xét tính đúng, sai của các khẳng định sau:
	\choiceTF
	{\True Điều kiện của ẩn là $x > 0, y > 0$.}
	{\True Hệ phương trình lập được từ bài toán là $\begin{cases} 3x + 2y = 36 \\ 2x + 4y = 40 \end{cases}$.}
	{Giá tiền của một chiếc bút là $8\,000$ đồng.}
	{\True Nếu mua $5$ quyển vở và $5$ chiếc bút cùng loại thì hết $70\,000$ đồng.}
	\loigiai{
		\begin{itemchoice}
			\itemch Đúng. Do $x, y$ là số tiền nên phải lớn hơn $0$.
			\itemch Đúng. Dựa vào giả thiết, ta lập được phương trình số tiền mua vở và bút.
			\itemch Sai. Giải hệ phương trình ta được $x=8$ (giá vở là $8\,000$ đ) và $y=6$ (giá bút là $6\,000$ đ).
			\itemch Đúng. Số tiền cần trả là $5 \times 8 + 5 \times 6 = 70$ (nghìn đồng).
		\end{itemchoice}
	}
\end{ex}

% nguồn: De_AI_Toan9_cuoi_hoc_ki_i_1 (nhập 18/11/2026)
\begin{ex}%[ma:d5e2b286]
	Cho biểu thức $A = \left( \frac{1}{\sqrt{x} - 1} + \frac{1}{\sqrt{x} + 1} \right) \cdot \frac{x - 1}{\sqrt{x}}$. Xét tính đúng, sai của các khẳng định sau:
	\choiceTF
	{Điều kiện xác định của biểu thức $A$ là $x \ge 0, x \ne 1$.}
	{\True Mẫu thức chung của biểu thức trong ngoặc là $x - 1$.}
	{Giá trị của biểu thức $A$ thay đổi phụ thuộc vào giá trị của biến $x$.}
	{\True Với $x = 4$, giá trị của $A$ bằng $2$.}
	\loigiai{
		\begin{itemchoice}
			\itemch Sai. Điều kiện xác định phải là $x > 0$ và $x \ne 1$ vì có $\sqrt{x}$ nằm ở mẫu.
			\itemch Đúng. Vì $(\sqrt{x}-1)(\sqrt{x}+1) = x-1$.
			\itemch Sai. Rút gọn $A = \frac{2\sqrt{x}}{x-1} \cdot \frac{x-1}{\sqrt{x}} = 2$. Giá trị $A$ luôn bằng $2$ với mọi $x$ thoả mãn ĐKXĐ.
			\itemch Đúng. Vì $A = 2$ với mọi $x > 0, x \ne 1$, nên tại $x=4$ thì $A=2$.
		\end{itemchoice}
	}
\end{ex}

% nguồn: De_AI_Toan9_cuoi_hoc_ki_i_1 (nhập 18/11/2026)
\begin{ex}%[ma:84a18a51]
	Một chiếc thang dài $5\text{ m}$ được tựa vào một bức tường thẳng đứng. Biết khoảng cách từ chân thang đến chân tường là $3\text{ m}$. Xét tính đúng, sai của các khẳng định sau:
	\choiceTF
	{\True Khoảng cách từ mặt đất lên đến điểm thang chạm tường là $4\text{ m}$.}
	{Góc tạo bởi chiếc thang và mặt đất (làm tròn đến độ) là $37^\circ$.}
	{\True Nếu đẩy chân thang vào gần tường sao cho chân thang cách tường $2\text{ m}$ thì điểm thang chạm tường cách mặt đất $\sqrt{21}\text{ m}$.}
	{\True Để thang chạm tường ở độ cao $4{,}8\text{ m}$ thì chân thang phải cách chân tường $1{,}4\text{ m}$.}
	\loigiai{
		\begin{itemchoice}
			\itemch Đúng. Dùng định lý Pythagore: $\sqrt{5^2 - 3^2} = 4\text{ (m)}$.
			\itemch Sai. Gọi góc đó là $\alpha$, ta có $\cos \alpha = \frac{3}{5} \Rightarrow \alpha \approx 53^\circ$.
			\itemch Đúng. Áp dụng định lý Pythagore lúc này: $\sqrt{5^2 - 2^2} = \sqrt{21}\text{ (m)}$.
			\itemch Đúng. Chân thang cách tường: $\sqrt{5^2 - 4{,}8^2} = \sqrt{25 - 23{,}04} = \sqrt{1{,}96} = 1{,}4\text{ (m)}$.
		\end{itemchoice}
	}
\end{ex}

% nguồn: De_AI_Toan9_cuoi_hoc_ki_i_1 (nhập 18/11/2026)
\begin{ex}%[ma:361a7aab]
	Cho hai tiếp tuyến $AB$ và $AC$ của đường tròn $(O)$ cắt nhau tại điểm $A$ nằm ngoài đường tròn ($B, C$ là các tiếp điểm). Xét tính đúng, sai của các khẳng định sau:
	\choiceTF
	{\True $AB = AC$.}
	{\True $OA$ là đường trung trực của đoạn thẳng $BC$.}
	{Góc $\widehat{BAO}$ luôn bằng góc $\widehat{AOC}$.}
	{\True Nếu $\widehat{BAC} = 60^\circ$ thì tam giác $ABC$ là tam giác đều.}
	\loigiai{
		\begin{itemchoice}
			\itemch Đúng. Theo tính chất hai tiếp tuyến cắt nhau.
			\itemch Đúng. Ta có $AB=AC$ và $OB=OC$ nên $A, O$ cùng nằm trên đường trung trực của $BC$.
			\itemch Sai. $AO$ là phân giác của $\widehat{BAC}$ nên $\widehat{BAO} = \widehat{CAO}$. Góc $\widehat{BAO}$ và $\widehat{AOC}$ là hai góc phụ nhau chứ không bằng nhau (trừ phi $\widehat{BAC}=90^\circ$).
			\itemch Đúng. Tam giác $ABC$ cân tại $A$ ($AB=AC$) có một góc bằng $60^\circ$ nên là tam giác đều.
		\end{itemchoice}
	}
\end{ex}

% nguồn: De_AI_Toan9_cuoi_hoc_ki_i_1 (nhập 18/11/2026)
\begin{bt}%[ma:39952af3]
	Tìm giá trị của $S = x + y$, biết $(x; y)$ là nghiệm của hệ phương trình $\begin{cases} 2x + y = 5 \\ 3x - 2y = 4 \end{cases}$.
	\shortans{3}
	\loigiai{
		Nhân phương trình đầu với $2$, ta có $\begin{cases} 4x + 2y = 10 \\ 3x - 2y = 4 \end{cases}$.
Cộng vế theo vế được $7x = 14 \Rightarrow x = 2$. \\
Thay vào phương trình đầu: $2(2) + y = 5 \Rightarrow y = 1$. \\
Vậy $S = x + y = 2 + 1 = 3$.
	}
\end{bt}

% nguồn: De_AI_Toan9_cuoi_hoc_ki_i_1 (nhập 18/11/2026)
\begin{bt}%[ma:69afa6b5]
	Tìm nghiệm nguyên lớn nhất của bất phương trình $\frac{2x - 1}{3} < \frac{x + 2}{2}$.
	\shortans{7}
	\loigiai{
		Quy đồng khử mẫu (nhân hai vế với $6 > 0$): $2(2x - 1) < 3(x + 2)$. \\
Khai triển: $4x - 2 < 3x + 6$. \\
Chuyển vế: $x < 8$. \\
Nghiệm nguyên lớn nhất nhỏ hơn $8$ là $7$.
	}
\end{bt}

% nguồn: De_AI_Toan9_cuoi_hoc_ki_i_1 (nhập 18/11/2026)
\begin{bt}%[ma:2d62835f]
	Tính giá trị của biểu thức $B = \sqrt{7 - 4\sqrt{3}} + \sqrt{3}$.
	\shortans{2}
	\loigiai{
		Ta có $7 - 4\sqrt{3} = 4 - 4\sqrt{3} + 3 = (2 - \sqrt{3})^2$. \\
Suy ra $\sqrt{7 - 4\sqrt{3}} = \sqrt{(2 - \sqrt{3})^2} = |2 - \sqrt{3}| = 2 - \sqrt{3}$ (vì $2 > \sqrt{3}$). \\
Khi đó $B = (2 - \sqrt{3}) + \sqrt{3} = 2$.
	}
\end{bt}

% nguồn: De_AI_Toan9_cuoi_hoc_ki_i_1 (nhập 18/11/2026)
\begin{bt}%[ma:81dee1d7]
	Một tia nắng mặt trời tạo với mặt đất một góc $45^\circ$. Cùng lúc đó, một tòa tháp đổ bóng trên mặt đất dài $25\text{ m}$. Hỏi chiều cao của tòa tháp là bao nhiêu mét?
	\shortans{25}
	\loigiai{
		Gọi chiều cao tòa tháp là $h$. \\
Tháp và bóng của nó trên mặt đất tạo thành hai cạnh góc vuông của một tam giác vuông. \\
Ta có: $\tan 45^\circ = \frac{\text{đối}}{\text{kề}} = \frac{h}{25}$. \\
Suy ra $h = 25 \cdot \tan 45^\circ = 25 \cdot 1 = 25\text{ (m)}$.
	}
\end{bt}

% nguồn: De_AI_Toan9_cuoi_hoc_ki_i_1 (nhập 18/11/2026)
\begin{bt}%[ma:676fd62e]
	Cho đường tròn $(O; 10\text{ cm})$ và một dây cung $AB = 16\text{ cm}$. Tính khoảng cách từ tâm $O$ đến dây cung $AB$ (đơn vị: cm).
	\shortans{6}
	\loigiai{
		Kẻ $OH \perp AB$ tại $H$. Theo quan hệ vuông góc giữa đường kính và dây, $H$ là trung điểm của $AB$. \\
Ta có $AH = \frac{AB}{2} = 8\text{ (cm)}$. \\
Áp dụng định lý Pythagore trong tam giác vuông $OHA$: $OH = \sqrt{OA^2 - AH^2} = \sqrt{10^2 - 8^2} = \sqrt{36} = 6\text{ (cm)}$.
	}
\end{bt}

% nguồn: De_AI_Toan9_cuoi_hoc_ki_i_1 (nhập 18/11/2026)
\begin{bt}%[ma:763c984e]
	Cho đường tròn $(O; 5\text{ cm})$ và một điểm $P$ nằm trong đường tròn sao cho $OP = 3\text{ cm}$. Qua $P$ kẻ dây cung $AB$. Độ dài nhỏ nhất của dây $AB$ bằng bao nhiêu cm?
	\shortans{8}
	\loigiai{
		Độ dài dây $AB$ nhỏ nhất khi khoảng cách từ $O$ đến dây $AB$ lớn nhất. \\
Kẻ $OH \perp AB$, ta luôn có $OH \le OP$. Vậy $OH$ lớn nhất bằng $OP = 3\text{ cm}$ (xảy ra khi $AB \perp OP$ tại $P$). \\
Khi đó, $AH = \sqrt{R^2 - OH^2} = \sqrt{5^2 - 3^2} = 4\text{ (cm)}$. \\
Độ dài dây nhỏ nhất là $AB = 2AH = 8\text{ (cm)}$.
	}
\end{bt}
